\section{Worum es geht}

\textbf{[SETZUNG]} Dieses Dokument entwickelt die vollständige Lagrange-Formulierung
der Gravitationsdynamik in FFGFT: das intrinsische Zeitfeld $\mathcal T=1/\max(m,\omega)$
als fundamentales Objekt, die Feldgleichung für sphärische und ausgedehnte Systeme,
die Kopplung aller Materiefelder an das Zeitfeld durch konforme Transformation, die
modifizierte Dirac-Gleichung und die vollständige Gesamt-Lagrange-Dichte. Das Dokument
liefert damit den Rahmen, der in A140 als „Gravitationsdynamik" benannt und dort nur
als Brücken-Ergebnis geführt ist. Was nach diesem Dokument noch offen bleibt, ist
explizit benannt.

\section{Das intrinsische Zeitfeld}

\textbf{[B]} Das fundamentale Objekt der Gravitationsdynamik in FFGFT ist nicht die
Metrik, sondern das intrinsische Zeitfeld:
\[
  \boxed{\;\mathcal T(x,t) = \frac{1}{\max(m,\omega)}\;}
\]
in natürlichen Einheiten ($\hbar=c=1$). $[\mathcal T]=[E^{-1}]$. Das Feld ist
reziprok zur lokalen Energie -- schwere Objekte haben kleine $\mathcal T$ (schnell
oscillierend), leichte haben große $\mathcal T$. Die Zeit-Masse-Dualität
$\tilde T\cdot m=1$ (A020/A060) ist die Nullfeldgleichung dieser Struktur.

\textbf{[B]} \textbf{Drei fundamentale Feldgeometrien:}

\begin{center}
\renewcommand{\arraystretch}{1.3}
{\small
\begin{tabular}{ll}
\toprule
Geometrie & Parameter \\
\midrule
Lokalisiert sphärisch &
  $\beta=2Gm/r$ (dim.los), $\xi=2\sqrt G\cdot m$ (dim.los),
  $\mathcal T(r)=(1/m_0)(1-\beta)$ \\
Lokalisiert nicht-sphärisch &
  $\beta_{ij}=r_{0ij}/r$ (Tensor), $\xi_{ij}=2\sqrt G\cdot I_{ij}$ \\
Unendlich homogen &
  $\nabla^2 m=4\pi G\rho_0 m+\Lambda_T m$, $\xi_\text{eff}=\xi/2$ \\
\bottomrule
\end{tabular}
}
\renewcommand{\arraystretch}{1.0}
\end{center}

\textbf{[K]} Für Quantencomputing-Anwendungen und atomare Skalen
($R/r\gg10^{12}$) ist nur die lokalisierte sphärische Geometrie relevant; die
anderen sind für theoretische Vollständigkeit.

\section{Die Feldgleichung}

\textbf{[B]} Für eine Punktmassenquelle $\rho=m\delta^3(\mathbf r)$ lautet die
vollständige geometrische Lösung:
\[
  m(r) = m_0\Bigl(1+\frac{2Gm}{r}\Bigr) = m_0(1+\beta),
  \qquad
  \mathcal T(r) = \frac{1}{m_0}(1+\beta)^{-1}\approx\frac{1}{m_0}(1-\beta).
\]
Der Faktor 2 in $r_0=2Gm$ folgt aus der relativistischen Feldstruktur und stimmt
exakt mit dem Schwarzschild-Radius überein. Die Feldgleichung:
\[
  \nabla^2 m = 4\pi G\,\rho(x,t)\cdot m.
  \qquad [E^3 = E^{-2}\cdot E^4\cdot E = E^3\;]\checkmark
\]
Für unendliche Systeme mit Hintergrund-Konsistenz:
\[
  \nabla^2 m = 4\pi G\rho_0 m + \Lambda_T m,
  \qquad \Lambda_T = -4\pi G\rho_0.
\]
Das $\Lambda_T$-Term erzeugt den kosmischen Abschirmeffekt $\xi_\text{eff}=\xi/2$.

\textbf{[B]} \textbf{Keine neuen Objekte.} Weder $r_0=2Gm$ noch
$\Lambda_T=-4\pi G\rho_0$ sind eigenständige Größen neben der Skalenstruktur
aus $\xi$. $r_0$ ist die charakteristische Länge \emph{der jeweiligen Masse} --
dieselbe Reziprozität wie $L_0$, nur mit $m$ statt der Planck-Masse als
Referenz. $\Lambda_T$ ist dieselbe Feldgleichung, gerechnet vom homogenen
Hintergrund statt von der Punktquelle aus. Beides ist Referenzpunktwahl,
nicht neue Physik -- dieselbe Logik wie die Referenzpunkt-Tabelle in A262.
Die verschiedenen Symbole dienen der Buchhaltung (welcher Referenzpunkt),
nicht einer ontologischen Trennung.

\section{Der universelle Skalenparameter aus dem Higgs-Matching}

\textbf{[K]} Der fundamentale Skalenparameter $\xi$ ist durch die Higgs-Physik
eindeutig bestimmt:
\[
  \boxed{\;\xi = \frac{\lambda_h^2\,v^2}{16\pi^3\,m_h^2}
  \approx 1{,}33\times10^{-4}\;}
\]
mit Higgs-Selbstkopplung $\lambda_h\approx0{,}13$, Higgs-VEV $v\approx246\,\text{GeV}$,
Higgs-Masse $m_h\approx125\,\text{GeV}$. Dimensionale Verifikation:
$[\xi]=[\lambda_h^2 v^2]/[16\pi^3 m_h^2]=[E^2]/[E^2]=[1]$. Die Bedingung
$\beta_T=1$ in natürlichen Einheiten ($\hbar=c=\alpha_\text{EM}=\beta_T=1$)
fixiert $\xi$ eindeutig und eliminiert alle freien Parameter.

\section{Die Lagrange-Formulierung}

\textbf{[B]} Die Zeitfeld-Lagrange-Dichte:
\[
  \mathcal L_\text{Zeit}
  = \sqrt{-g}\,\Bigl[
    \tfrac12\,g^{\mu\nu}\partial_\mu\mathcal T\,\partial_\nu\mathcal T
    - V(\mathcal T)
  \Bigr].
\]
Dimensionale Konsistenz: $[\mathcal L_\text{Zeit}]=[E^{-4}]([E^2]+[E^4])=[E^0]$.

\textbf{[B]} \textbf{Materiefeld-Kopplung durch konforme Transformation.} Alle
Materiefelder koppeln an das Zeitfeld über:
\[
  g_{\mu\nu}\;\to\;\Omega^2(\mathcal T)\,g_{\mu\nu},
  \qquad\Omega(\mathcal T)=\frac{\mathcal T_0}{\mathcal T}\quad[\text{dim.los}].
\]

Skalare Felder:
\[
  \mathcal L_\phi
  = \sqrt{-g}\,\Omega^4(\mathcal T)\,
    \Bigl(\tfrac12 g^{\mu\nu}\partial_\mu\phi\,\partial_\nu\phi
    - \tfrac12 m^2\phi^2\Bigr).
  \qquad [\mathcal L_\phi]=[E^{-4}][1][E^4]=[E^0]\checkmark
\]

Fermion-Felder:
\[
  \mathcal L_\psi
  = \sqrt{-g}\,\Omega^4(\mathcal T)\,
    \bigl(i\bar\psi\gamma^\mu\partial_\mu\psi - m\bar\psi\psi\bigr).
  \qquad [\mathcal L_\psi]=[E^{-4}][1][E^4]=[E^0]\checkmark
\]

Higgs-Zeitfeld-Kopplung:
\[
  \mathcal L_{\text{Higgs-T}}
  = |\mathcal D_{\text{Higgs-T}}|^2 - V(\mathcal T,\Phi),
  \qquad
  \mathcal D_{\text{Higgs-T}}
  = \mathcal T(\partial_\mu+igA_\mu)\Phi + \Phi\,\partial_\mu\mathcal T.
\]
Das liefert $\mathcal T=1/(y\langle\Phi\rangle)$ -- die Verbindung zwischen
Zeitfeld und Higgs-Vakuumerwartungswert.

\textbf{[B]} \textbf{Vollständige Gesamt-Lagrange-Dichte:}
\[
  \mathcal L_\text{Gesamt}
  = \mathcal L_\text{Zeit}
  + \mathcal L_\text{Eich}
  + \mathcal L_\phi
  + \mathcal L_\psi
  + \mathcal L_{\text{Higgs-T}},
\]
mit Eich-Term $\mathcal L_\text{Eich}=\sqrt{-g}\,(-\tfrac14 F_{\mu\nu}F^{\mu\nu})$.
Jeder Summand ist dimensional $[E^0]$.

\section{Die modifizierte Dirac-Gleichung}

\textbf{[B]} Im FFGFT-Zeitfeld-Framework lautet die Dirac-Gleichung:
\[
  \bigl[i\gamma^\mu(\partial_\mu+\Gamma_\mu^{(T)})-m(x,t)\bigr]\psi=0,
\]
mit der Zeitfeld-Verbindung:
\[
  \Gamma_\mu^{(T)}
  = \frac{1}{\mathcal T}\partial_\mu\mathcal T
  = -\frac{\partial_\mu m}{m^2}.
\]
Diese Gleichung reduziert sich im Grenzfall $\mathcal T\to\text{const}$
($\Gamma_\mu^{(T)}\to0$) auf die Standard-Dirac-Gleichung. Die Verbindung zur
Schrödinger-Gleichung (A075) bleibt über die Polarzerlegung des Vakuumfelds.

\textbf{[B]} \textbf{Modifizierte Schrödinger-Gleichung:}
\[
  i\,\mathcal T\,\frac{\partial\Psi}{\partial t}
  + i\,\Psi\,\Bigl[\frac{\partial\mathcal T}{\partial t}
  + \mathbf v\cdot\nabla\mathcal T\Bigr]
  = \hat H\,\Psi.
\]
Im Grenzfall $\partial_\mu\mathcal T\to0$ reduziert sie sich auf die
Standardform $i\hbar\,\partial_t\Psi=\hat H\Psi$ mit $\mathcal T=1/m\leftrightarrow\hbar$.

\section{Dimensionale Konsistenz-Überblick}

\textbf{[K]}
\begin{center}
\renewcommand{\arraystretch}{1.3}
{\small
\begin{tabular}{lccl}
\toprule
Gleichung & Linke Seite & Rechte Seite & Status \\
\midrule
Zeitfeld-Definition & $[E^{-1}]$ & $[1/\max(m,\omega)]=[E^{-1}]$ & \checkmark \\
Feldgleichung & $[E^3]$ & $[4\pi G\rho m]=[E^3]$ & \checkmark \\
$\beta$-Parameter & $[1]$ & $[2Gm/r]=[1]$ & \checkmark \\
$\xi$ aus Higgs & $[1]$ & $[\lambda_h^2 v^2/(16\pi^3 m_h^2)]=[1]$ & \checkmark \\
Lagrange-Dichte & $[E^0]$ & alle Terme & \checkmark \\
\bottomrule
\end{tabular}
}
\renewcommand{\arraystretch}{1.0}
\end{center}

\section{Testbare Vorhersagen}

\textbf{[K]} Der Lagrange-Rahmen liefert spezifische Vorhersagen:
\[
  a_e^{(T0)} = \frac\alpha{2\pi}\cdot\xi^2\cdot I_\text{Schleife}
  \approx 2{,}34\times10^{-10}
\]
(QED-Korrektur zum anomalen magnetischen Moment). Rotationskurven mit
modifiziertem Potential: $v^2(r)=GM/r+\kappa r^2$. Energieabhängige
Zeitverzögerung: $\Delta t=(\xi/c)(1/E_1-1/E_2)(2Gm/r)$.

\section{Prüfskript}

\begin{itemize}[leftmargin=2em,itemsep=2pt,topsep=2pt]
  \item Dimensionale Konsistenz aller Terme numerisch; Grenzfall Dirac→Standard;
    $\xi$ aus Higgs-Matching:
    \texttt{python/A\_Serie\_Skripte/a142\_gravitation\_lagrange.py}
\end{itemize}


\section{Zeichenerklärung}

Symbole die in der gesamten A-Serie verwendet werden, sind in A010 erklärt.
Spezifisch für dieses Dokument:

\begin{center}
\renewcommand{\arraystretch}{1.3}
{\small\begin{tabular}{lp{9cm}}
\toprule
Symbol & Bedeutung \\
\midrule
$\mathcal{T}(x,t)$ & Intrinsisches Zeitfeld (dynamisches Skalarfeld) \\
$\Omega$ & Konforme Transformationsmatrix \\
$\Gamma_\mu^{(T)}=\partial_\mu m/m^2$ & Zeitfeld-Zusammenhang (Verbindungskoeffizient) \\
$\mathcal{L}_\text{Zeit}$ & Lagrange-Dichte des Zeitfelds \\
$G_{\mu\nu}$ & Einstein-Tensor \\
$T_{\mu\nu}$ & Energie-Impuls-Tensor \\
\bottomrule
\end{tabular}}
\renewcommand{\arraystretch}{1.0}
\end{center}

\section{Was offen bleibt}

\textbf{Eine separate Quantisierung der Gravitation ist in FFGFT nicht
gefordert.} $G$ ist keine unabhängige Feldgröße, sondern folgt intrinsisch
aus $\xi$: $G=\xi^2/(4m_\text{char})$. Gravitation ist die Projektion der
Zeitfeld-Geometrie, kein eigenständiges Quantenfeld. Die Frage
„wann wird Gravitation quantisiert?" stellt sich in diesem Rahmen nicht --
sie ist durch die Grundstruktur bereits beantwortet. Was bleibt: die
Überführung der klassischen Feldgleichungen in explizite Schleifenkorrekturen
für Szenarien jenseits des Schwarzschild-Grenzfalls.

\textbf{Der Übergang zu messbaren Vorhersagen jenseits des Schwarzschild-Falls
ist nicht vollständig.} Die QED-Korrektur (oben) ist berechenbar; Vorhersagen
für starke Felder, Gravitationswellen oder den post-newtonschen Sektor sind nicht
systematisch ausgeführt.

\textbf{Die Verbindung zur Projektionskette (A090) fehlt.} Wie der
Gravitationssektor sich in die Typ-I/II/III-Struktur einfügt -- ob er ein
Typ-II-Verlust oder ein Typ-I-Äquivalent ist -- ist nicht expliziert.

\section{Ersetzt}

Dieses Dokument nimmt auf: Dok.~067 (Zeitfeld-Lagrange-Dichte mit dimensionaler
Konsistenz, Materiefeld-Kopplung durch konforme Transformationen, sphärisch
symmetrische Lösungen, Higgs-Kopplung, $\xi$ aus Higgs-Matching, vollständige
Gesamt-Lagrange-Dichte, modifizierte Dirac-Gleichung, experimentelle Vorhersagen).
\emph{Nicht} übernommen: der Abschnitt zu einem statischen Universum aus Dok.~067
-- er gehört nicht in den A-Serien-Kern (A010).

\section{Verwendung von KI-Werkzeugen}

Dieses Dokument ist unter Verwendung KI-gestützter Werkzeuge entstanden. Die
Fragestellungen, die Setzungen und alle inhaltlichen Entscheidungen stammen vom
Autor; die Werkzeuge formulieren aus, rechnen nach und erstellen Prüfskripte.
Die Verantwortung für Inhalt und Fehler liegt vollständig beim Autor. Der
ausführliche Wortlaut steht in A010.
